% Section 18.4: The structural meaning: $\dd$ is the adjoint of $\partial$ | Part: 3 | Depends on: sec:intk, sec:fields | Status: todo | Original lines: 2755-2768
%!TEX root = ../../../main.tex
\subsection{The structural meaning: $\dd$ is the adjoint of $\partial$}\label{sec:adjoint}

The modern reading of \eqref{eq:stokes} is obtained by introducing the pairing between $k$-forms and $k$-chains $c$ (formal sums of oriented $k$-dimensional pieces): $\langle c,\alpha\rangle:=\int_c\alpha$. Then Stokes' theorem says
\begin{equation}\label{eq:adjoint}
\langle c,\dd\alpha\rangle=\langle\partial c,\alpha\rangle,
\end{equation}
\ie the exterior derivative and the boundary operator are \emph{adjoint} to each other. Consequences:
\begin{itemize}
\item $\partial^2=0$ (the boundary of a boundary is empty) is the dual of $\dd^2=0$, and each explains the other. Indeed, applying \eqref{eq:adjoint} twice: $\langle c,\dd\dd\alpha\rangle=\langle\partial c,\dd\alpha\rangle=\langle\partial\partial c,\alpha\rangle$. If $\partial\partial c=0$ the right-hand side vanishes for every $c$, and a form whose integral over every small piece is zero must be zero: $\dd\dd\alpha=0$.
\item Integration is a chain map from the de Rham complex $(\Omega^\bullet(M),\dd)$ to the singular cochain complex, hence it descends to cohomology. de Rham's theorem (1931): this map $H^k_{\mathrm{dR}}(M)\to H^k(M;\R)$ is an isomorphism. Closed forms integrate to the same value over homologous cycles (``periods''), as in Example~\ref{ex:angle}.
\item Stokes' theorem is \emph{metric-free} and uses only a smooth structure and an orientation. Metrics enter only when we translate forms into vector fields.
\item Applications reach well beyond calculus: degree theory and the Brouwer fixed-point theorem, the Poincar\'e lemma and Poincar\'e duality, conservation laws, and Maxwell's equations $\dd F=0$, $\dd\hodge F=J$, whose integral forms are the laws of Faraday and Gauss.
\end{itemize}
